\section{An obstruction for essentially self-dual spin cores}
\label{sec:spin-core-obstruction}

The sporadic calculation does not cover alternating simple factors.
The following separate argument removes the essentially self-dual
spin case in all alternating degrees at least eight.  It does not
remove ordinary alternating representations or nonselfdual spin
representations.

\begin{lemma}[HD-SPIN-01a: the spin-normalizer moving rank]
\label{HD-SPIN-01a}
Let $N=2.\mathfrak A_n$, $n\ge8$, act faithfully and irreducibly
on a complex space $W$ of dimension $a$.  Every nonscalar
finite-order element $g$ of its linear normalizer satisfies
\[
                  \nu(g)\ge a/8,                         \tag{SPIN.1}
\]
where $\nu$ is the codimension of a largest eigenspace.
The same conclusion holds if $g$ has finite order only after a
scalar rescaling.
\end{lemma}

\begin{proof}
We use the local character-table input appearing in the complete
proof of \cite[Lemma~2.23, pp.~621--622]{HD:GT2012}: every faithful
irreducible character $\psi$ of $2.\mathfrak A_6$ satisfies
\[
       |\psi(h)|\le\psi(1)/2\quad
       (h\notin Z(2.\mathfrak A_6)).                       \tag{SPIN.2}
\]
This input is also independently checked on the complete retained
table in the certificate described below.  We use the input (SPIN.2),
not the weaker final $7/8$ character-ratio conclusion of that lemma.

Schur's lemma identifies the linear centralizer of $N$ with the
scalars.  An automorphism of the perfect group $N$ inducing the
identity on $\mathfrak A_n$ has the form $x\mapsto z(x)x$, with
$z$ a homomorphism to the center, and is therefore the identity.
Thus the normalizer modulo scalars embeds in
$\operatorname{Aut}(\mathfrak A_n)=\mathfrak S_n$ for $n\ge8$;
the latter standard automorphism description is also used in
\cite[proof of Proposition~3.4, p.~626]{HD:GT2012}.
Let $\pi\ne1$ be the permutation induced by $g$.

There is a three-cycle $c$ not centralized by $\pi$.  Indeed, the
three-cycles generate $\mathfrak A_n$, whose centralizer in its
natural permutation group is trivial.  The last assertion can
also be checked directly: if $\pi$ commutes with every three-cycle,
it preserves the support of every three-cycle, hence every
three-element subset; the intersections of these subsets recover
each singleton for $n\ge4$.  Choose a lift $t\in N$ of $c$.
Then
\[
                   h=[g,t]\in N
\]
projects to the nonidentity product of the three-cycles
$\pi c\pi^{-1}$ and $c^{-1}$.  Its support is contained in at
most six letters.  It therefore belongs to a natural
$2.\mathfrak A_6$ subgroup of $N$ and is noncentral there.
This is precisely the natural-subgroup step in the proof of
\cite[Lemma~2.23]{HD:GT2012}; restriction of the standard double
cover to these six letters is the standard double cover.

The central involution of this subgroup acts as $-1$ on $W$.
Consequently every irreducible constituent of the restriction
to it is faithful on the center, and hence faithful, since a
proper normal subgroup of a quasisimple group is central.
Summing (SPIN.2) over the constituents gives
$|\operatorname{tr}_W(h)|\le a/2$.  If $m$ is the largest
eigenvalue multiplicity of $h$, rotate that eigenvalue to one.
The triangle inequality gives
\[
           2m-a\le|\operatorname{tr}_W(h)|\le a/2,
           \qquad \nu(h)=a-m\ge a/4.                     \tag{SPIN.3}
\]

Choose a scalar $\lambda$ with
$\operatorname{rank}(1-\lambda g)=\nu(g)$.
For any invertible matrices $A,B$,
$\operatorname{rank}(1-AB)\le
\operatorname{rank}(1-A)+\operatorname{rank}(1-B)$.
The commutator of $\lambda g$ with $t$ is still $h$, so
\[
 \nu(h)\le\operatorname{rank}(1-h)
 \le 2\operatorname{rank}(1-\lambda g)=2\nu(g).
\]
Together with (SPIN.3) this proves (SPIN.1).
\end{proof}

\begin{theorem}[HD-SPIN-01b: exclusion of essentially self-dual spin cores]
\label{HD-SPIN-01b}
Suppose a finite-cover Weil sixfold source has nonsplit actual
polarization, exact imaginary quadratic character field, and
period rank nine.  Let an irreducible primitive inducing block
have finite image $J$ and degree $e\ge9$, with arbitrary inducing
character field.  If $J$ is essentially self-dual, it cannot have
an irreducibly acting normal central product
\[
                  E=E_1\cdots E_t,\qquad
                  E_i\simeq2.\mathfrak A_{n_i},
                  \quad n_i\ge8,                        \tag{SPIN.4}
\]
whose factor representations are faithful.
In particular this removes the essentially self-dual spin-core
case from D-BI-19.  Permutations of isomorphic factors and all
compatible outer automorphisms are allowed.
\end{theorem}

\begin{proof}
Restriction to the commuting quasisimple factors gives
$V=\bigotimes_i V_i$, with $a_i=\dim V_i$.  The published
spin-degree lower bound
\[
                   a_i\ge2^{\lfloor n_i/2\rfloor-1}\ge8  \tag{SPIN.5}
\]
is recorded and used in
\cite[proof of Proposition~3.4, p.~626]{HD:GT2012}.
It concerns faithful complex spin characters of the double
alternating cover, not ordinary representations of $\mathfrak A_n$.

Every nonscalar element of $J$ has $\nu(g)\ge e/8$.
For an element preserving each tensor factor, this follows from
HD-SPIN-01a on a nonscalar factor and the tensor eigenspace bound:
a largest eigenspace multiplicity in a tensor product is at most
that in one factor times the remaining dimensions.  The factor
maps have finite order up to scalars, as follows by taking a power
of the tensor operator.
For an element permuting factors, the weighted cyclic-word
argument of D-SPOR-03, equation (SPOR.6), instead gives
$\nu(g)\ge e/4$.  That argument uses only the dimensions of
the factors, not any sporadic group hypothesis.  The intrinsic
factor-algebra normalizer argument in the same proof applies
unchanged to (SPIN.4).  Thus the weaker bound $e/8$ holds for
every nonscalar element.  Every genuine scalar inertia moves all
$e$ dimensions.

The actual intermediate-cover argument of D-BI-09 gives
$r\ge12$ genuine inertia elements with total moving rank
$2e+6$.  The essentially self-dual similitude character is
nontrivial because the inducing character field contains the
original imaginary quadratic field.  At least two inertias have
nontrivial multiplier, and each moves at least $\lceil e/2\rceil$
dimensions, by the isotropic-fixed-space proof in D-BI-09.
The other at least ten inertias move at least $\lceil e/8\rceil$.
Therefore
\[
       2\lceil e/2\rceil+10\lceil e/8\rceil\le2e+6,
       \qquad e\le24.                                   \tag{SPIN.6}
\]
The second inequality follows since the left side is at least
$9e/4$.  D-BI-09, with its complete primitive and tensor cases,
gives $e\ge24$.  Hence $e=24$.

There can be only one factor in (SPIN.4), since two factors
already give $e\ge8^2=64$.  Equation (SPIN.5) also forces
$8\le n\le11$.  The complete faithful-on-center character rows
of the published tables $2.A8,2.A9,2.A10,2.A11$ have the
following degree-twenty-four data:
\[
\begin{array}{c|c|c}
\text{table}&\text{degree-24 rows}&\text{self-duality}\\ \hline
2.A8&16,\ 17&\text{both nonreal}\\
2.A9&\varnothing&---\\
2.A10&\varnothing&---\\
2.A11&\varnothing&---
\end{array}                                             \tag{SPIN.7}
\]
Here rows are the exact CTblLib identifiers in the retained
certificate, not an assumption that all spin degrees are powers
of two.  In $2.A8$ the two rows are complex conjugates, each
with Frobenius--Schur indicator zero.  Their values include
$b_7=\zeta_7+\zeta_7^2+\zeta_7^4$ and $\overline b_7$;
the character field is $\mathbf Q(\sqrt{-7})$.

Finally, the multiplier of an essentially self-dual $J$-module
is a one-dimensional character.  Its restriction to the perfect
group $E$ is trivial.  The invariant nondegenerate bilinear form
would therefore make the irreducible $E$-module self-dual.
Both actual degree-twenty-four rows in (SPIN.7) contradict this.
This finishes the exclusion.
\end{proof}

\paragraph{Certificate and scope.}
The complete five small-cover tables, their source strings and
power maps are retained alongside the alternating-cover
certificate in \texttt{agents/imprimitive/sporadic/}.
They are published ATLAS inputs accessed through the exact
identifiers of \cite{SPOR:CTblLib}.  The numerical verification
checks (SPIN.2), the full faithful degree lists, the degree-24
exception and its nonreal witnesses; it does not claim to prove
the general spin-degree lower bound or the identification of
the tables with the named groups.  The latter premises retain
their explicit published locators above.
The theorem leaves nonselfdual spin cores, all ordinary
alternating cores, and the small covers in degrees $n<8$
outside its scope.
